
%-------------------------------------------
\paragraph{Strain rate and rotation rate}

The strain rate (corresponding to the velocity field $\vec\upnu$) is given by 
\[
\dot{\bm \varepsilon} = \frac12 ( \vec\nabla \vec\upnu + \vec\nabla \vec\upnu^T) 
\]
while the corresponding vorticity tensor
(sometimes called the 'spin tensor', e.g. \cite{moql07}, or 'vorticity tensor')  is given by
\[
\dot{\bm \omega} = \frac12 \left( \vec\nabla \vec\upnu - \vec\nabla \vec\upnu^T \right) 
\]
Obviously we find that $\dot{\bm \varepsilon}+\dot{\bm \omega}=\vec\nabla \vec\upnu  = {\bm L}$.



%---------------------------
\paragraph{Viscous rheology}

From MD syllabus:
\begin{eqnarray}
{\bm \sigma} 
&=& -p {\bm 1} + {\bm \pi} \nn\\
&=& -p {\bm 1} + \lambda {\rm tr}(\dot{\bm \varepsilon}) {\bm 1} + 2 \eta \dot{\bm \varepsilon} \nn\\
&=& -p {\bm 1} + \xi {\rm tr}(\dot{\bm \varepsilon}) {\bm 1} + 2 \eta \dot{\bm \varepsilon}^d 
\end{eqnarray}
where $p$ is the thermodynamic pressure and 
where $\xi=\lambda+2\eta/3$ is the bulk viscosity. It follows that 
\begin{equation}
{\rm tr}({\bm \sigma}) = -3p + 3 \xi {\rm tr}(\dot{\bm \varepsilon}) \label{eq:trsigmav}
\end{equation}
and automatically the deviatoric stress ${\bm \tau}$ follows:
\begin{equation}
{\bm \tau} = {\bm \sigma} - \frac13 {\rm tr}({\bm \sigma}) {\bm 1} = 2 \eta \dot{\bm \varepsilon}^d \label{eq:ve:tauv}
\end{equation}


%---------------------------
\paragraph{Elastic rheology}

\begin{equation}
{\bm \sigma} = \lambda {\rm tr}({\bm \varepsilon}) {\bm 1} + 2 \mu {\bm \varepsilon}
\end{equation}
where $\lambda$ is the Lame parameter and $\mu$ the shear modulus.

Then 
\begin{equation}
{\rm tr}({\bm \sigma}) 
= 3 \lambda {\rm tr}({\bm \varepsilon}) + 2 \mu {\rm tr}({\bm \varepsilon}) 
= 3 \left(\lambda + \frac{2\mu}{3}\right) {\rm tr}({\bm \varepsilon})
= 3 K {\rm tr}({\bm \varepsilon}) \label{eq:trsigmae}
\end{equation}
where $K$ is the elastic bulk modulus.

Like wise 
\begin{equation}
{\bm \tau} 
= {\bm \sigma} - \frac13 {\rm tr}({\bm \sigma}) {\bm 1}
= \lambda {\rm tr}({\bm \varepsilon}) {\bm 1} + 2 \mu {\bm \varepsilon}
- \frac13 \left(    3 \lambda {\rm tr}({\bm \varepsilon}) + 2 \mu {\rm tr}({\bm \varepsilon})    \right) {\rm tr}({\bm \varepsilon}) {\bm 1}
= 2 \mu {\bm \varepsilon}^d\label{eq:ve:taue}
\end{equation}

%---------------------------
\paragraph{total strain rate}

The total strain rate tensor (i.e. {\it not} the deviatoric one) is the sum of the viscous and the elastic 
strain rates
\begin{equation}
\dot{\bm \varepsilon}_{\rm T} = \dot{\bm \varepsilon}_{\rm v} + \dot{\bm \varepsilon}_{\rm e}
\label{eq:ve:epsT} 
\end{equation}

We have
\[
\dot{\bm\varepsilon}_{\rm e}^d = 
\dot{\bm\varepsilon}_{\rm e}  - \frac13 {\rm tr} (\dot{\bm\varepsilon}_{\rm e}) {\bm 1} 
\]
\[
\dot{\bm\varepsilon}_{\rm v}^d = 
\dot{\bm\varepsilon}_{\rm v}  - \frac13 {\rm tr} (\dot{\bm\varepsilon}_{\rm v}) {\bm 1} 
\]

We can then proceed to take the deviatoric part of the this equation: 
\begin{eqnarray}
\dot{\bm \varepsilon}_{\rm T}^d 
&=& 
\dot{\bm \varepsilon}_{\rm T} - \frac13 {\rm tr}(\dot{\bm \varepsilon}_{\rm T}) {\bm 1} \nn\\
&=& \dot{\bm \varepsilon}_{\rm v} + \dot{\bm \varepsilon}_{\rm e} 
- \frac13 {\rm tr}( \dot{\bm \varepsilon}_{\rm v} + \dot{\bm \varepsilon}_{\rm e} ) {\bm 1} \nn\\
&=& \dot{\bm \varepsilon}_{\rm v}^d + \dot{\bm \varepsilon}_{\rm e}^d  \label{eq:ve:epsTdev}
\end{eqnarray}


From Eq.~\eqref{eq:ve:tauv} and Eq.~\eqref{eq:ve:taue} 
(note that we must take its time derivative, 
which yields $\check{\bm\tau}$... I still need to fully wrap my head around this. 
See next section on this topic) we can write Eq.~\eqref{eq:ve:epsTdev}

\begin{equation}
\dot{\bm \varepsilon}_{\rm T}^d = 
\frac{1}{2\eta} {\bm \tau} + \frac{1}{2\mu} \check{\bm\tau}
\label{eq:ve:decomp}
\end{equation}
which forms the basis for all publications that implement elasticity in a viscous code.

Let us now turn to the volumetric part of Eq.~\eqref{eq:ve:epsT}: 
\begin{equation}
{\rm tr}(\dot{\bm \varepsilon}_{\rm T}) = {\rm tr}(\dot{\bm \varepsilon}_{\rm v}) + {\rm tr}(\dot{\bm \varepsilon}_{\rm e})
\end{equation}
Obviously if any or both deformation mechanisms are assumed to be incompressible this 
equation is not needed since all terms are zero. Also
it seems to be quite common to only allow for compressible elastic deformation while
the viscous one is incompressible (see for example Simpson's book, eq 12.43).

From Eq.~\eqref{eq:trsigmae} we obtain (again I write $\check{\bm \sigma}$): 
\begin{equation}
{\rm tr}(\dot{\bm \varepsilon}_e) = \frac{1}{3K} {\rm tr}(\check{\bm \sigma})
\end{equation}
Remark: In order to recover the $\dot{p}$ term of the papers, I need to state that $p=-\frac13 {\rm tr}({\bm \sigma})$, 
state that $\check{p}=\dot{p}$ because $p$ is a scalar (as in \cite{modm01}) and then I obtain:
\begin{equation}
{\rm tr}(\dot{\bm \varepsilon}_e) = - \frac{1}{K} \dot{p} 
\end{equation}


Conversely, from Eq.~\eqref{eq:trsigmav} we obtain:
\[
{\rm tr}(\dot{\bm \varepsilon}_{\rm v})  = ... 
\]

{\color{red} I am struggling here ... I need to get to :}


In the end in \cite{modm01} they arrive at
\[
\frac{\dot{p}}{K_{\rm e}}+\frac{p}{\xi} = - {\rm tr} (\dot{\bm \varepsilon})
\]
where K bulk modulus and $\xi$ bulk viscosity
Also in \cite{modm02}:
\[
\frac{1}{K} \frac{Dp}{Dt} + \frac{p}{\xi} = -{\rm tr}(\dot{\bm\varepsilon})
\]

{\color{red} I could arrive at these is I assume that the thermodynamic pressure $p$ is
to be discarded , so that the pressure above is in fact the 'mechanical pressure' (?). see section 18.2.3 of MD}

Ask Riad?


%------------------------------------------
\paragraph{About the weird time derivative}



\[
\check{\bm\tau} = 
\frac{{\cal D}{\bm \tau}}{{\cal D} t} = \frac{\partial{\bm \tau}}{\partial t}
+ \underbrace{(\vec\upnu \cdot \vec\nabla) {\bm \tau} }_{\rm advection}
+ \underbrace{{\bm \tau}\cdot \dot{\bm \omega} - \dot{\bm \omega}\cdot {\bm \tau}}_{\rm rotation}
\underbrace{- a ( {\bm \tau} \cdot \dot{\bm \varepsilon} + \dot{\bm \varepsilon} \cdot {\bm \tau})}_{\rm stretching}
\]
The labelling of the terms is to be found in \cite{moql07}.
${\cal D}/{\cal D} t$ is called an ``invariant time derivative'' in \cite{hard91}.

Sometimes the material derivative $\frac{D}{Dt} = \frac{\partial }{\partial t} + \vec\upnu \cdot \vec\nabla$
is used so that the equation above then becomes:
\[
\check{\bm\sigma} = 
\frac{{\cal D}{\bm \sigma}}{{\cal D} t} = \frac{D {\bm\tau}}{Dt}
+ {\bm \sigma}\cdot \dot{\bm \omega} - \dot{\bm \omega}\cdot {\bm \sigma}
- a ( {\bm \sigma} \cdot \dot{\bm \varepsilon} + \dot{\bm \varepsilon} \cdot {\bm \sigma})
\]

Note that in the equation above a minus sign is present in front of $a$ (which we 
find in \cite{hard91,momy93,daws16}) but we find a plus sign in \cite{moql07} (probably a typo?). 

In \cite{modm01} we read that ``the ${\bm\tau}\cdot \dot{\bm\omega}-\dot{\bm\omega}\cdot {\bm\tau}$ terms account 
for material spin during advection which reorients the elastic stored-stress tensor.''

In \cite{daws16} we read that 
in the case of large deformations, the material derivative of the Cauchy
stress tensor takes the form of the objective Gordon-Schowalter derivative (Saramito, 2016):
\[
\frac{{\cal D}{\bm \sigma}}{{\cal D} t} = \frac{\partial{\bm \sigma}}{\partial t}
+ (\vec\upnu \cdot \vec\nabla) {\bm \sigma} +
\beta_a (\vec\nabla\vec\upnu,{\bm \sigma})
\]
where the additional term $\beta_a$  accounts for the effects
of rotation and deformation on the evolution of the stress tensor 
and is expressed as 
\[
\beta_a (\vec\nabla\vec\upnu,{\bm \sigma})
= 
{\bm \sigma}\cdot \dot{\bm \omega} - \dot{\bm \omega}\cdot {\bm \sigma}
- a ( {\bm \sigma} \cdot \dot{\bm \varepsilon}
+ \dot{\bm \varepsilon} \cdot {\bm \sigma})
\]
where $-1 \le a \le +1$ is a free parameter. 
With the choice $a = 1$ the upper convected derivative is obtained, 
with $a = 0$ the co-rotational one \cite{hard91}.

In \cite{moql07} we read: ``When $a=0$, $\check{\bm \tau}$ is known as the Jaumann derivative, and
when $a=1$ or $a=-1$ it is known as the upper- or lower-convected Maxwell derivative, respectively. 
The choice $a=0$ considerably simplifies the formulation as it ensures 
that $\check{\bm \tau}$ is deviatoric when ${\bm \tau}$ is deviatoric.''

Note that in the appendix D of \cite{bepo10} the term ${\bm \sigma}\cdot \dot{\bm \omega} - \dot{\bm \omega}\cdot {\bm \sigma}$
is further worked out in 2d Cartesian coordinates and yields:
\[
{\bm \sigma}\cdot \dot{\bm \omega} - \dot{\bm \omega}\cdot {\bm \sigma}
=
\left(
\begin{array}{cc}
2 \omega \tau_{xy} & \omega (\tau_{yy}-\tau_{xx}) \\
\omega (\tau_{yy}-\tau_{xx}) -2 \omega \tau_{xy}
\end{array}
\right)
\]

%---------------------------------
\paragraph{Discretisation}

Let us for now drop the last term (i.e. $a=0$) so that 
\[
\check{\bm\tau} = 
\frac{{\cal D}{\bm \tau}}{{\cal D} t} = 
\frac{D {\bm \tau}}{D t}
+ {\bm \tau}\cdot \dot{\bm \omega} - \dot{\bm \omega}\cdot {\bm \tau}
\]
Following \cite{modm02} we discretise this as 
\[
\check{\bm\tau}^{t+\delta t} \simeq 
\frac{ {\bm\tau}^{t+\delta t}- {\bm\tau}^t   }{\delta t}
+ {\bm \tau}^t\cdot \dot{\bm \omega}^t - \dot{\bm \omega}^t \cdot {\bm \tau}^t
\]



Then Eq.~\eqref{eq:ve:decomp} is written at time $t+\delta t$:
\begin{eqnarray}
\dot{\bm \varepsilon}^{d,t+\delta t} 
&=& 
\frac{1}{2\eta} {\bm \tau}^{t+\delta t} + \frac{1}{2\mu} \check{\bm \tau}^{t+\delta t}  
\nn\\
\dot{\bm \varepsilon}^{d,t+\delta t} 
&=&
\frac{1}{2\eta} {\bm \tau}^{t+\delta t} + \frac{1}{2\mu} 
\left(
\frac{ {\bm\tau}^{t+\delta t}- {\bm\tau}^t   }{\delta t}
+ {\bm \tau}^t\cdot \dot{\bm \omega}^t - \dot{\bm \omega}^t \cdot {\bm \tau}^t
\right) 
\nn\\
\dot{\bm \varepsilon}^{d,t+\delta t} 
&=&
\left(\frac{1}{2\eta} + \frac{1}{2 \mu\delta t} \right){\bm \tau}^{t+\delta t} 
- \frac{1}{2 \mu\delta t}  {\bm\tau}^t 
+ \frac{1}{2 \mu} 
( {\bm \tau}^t\cdot \dot{\bm \omega}^t - \dot{\bm \omega}^t \cdot {\bm \tau}^t) \nn
\end{eqnarray}

or, 
\[
\frac12
\left(\frac{1}{\eta} + \frac{1}{\mu\delta t} \right){\bm \tau}^{t+\delta t} 
= \dot{\bm \varepsilon}^{d,t+\delta t} 
+ \frac{1}{2 \mu\delta t}  {\bm\tau}^t
- \frac{1}{2 \mu} ( {\bm \tau}^t\cdot \dot{\bm \omega}^t - \dot{\bm \omega}^t \cdot {\bm \tau}^t)
\]

Defining the visco-elastic viscosity as  
\[
\eta_{\rm ve} 
= \left( \frac{1}{\eta} + \frac{1}{\mu \delta t} \right)^{-1}
= \frac{\eta\mu \; \delta t}{\eta + \mu \delta t}
\]
we can write
\[
\frac{1}{2\eta_{\rm ve}} {\bm \tau}^{t+\delta t} 
= \dot{\bm \varepsilon}^{d,t+\delta t} 
+ \frac{1}{2 \mu\delta t}  {\bm\tau}^t
- \frac{1}{2 \mu} ( {\bm \tau}^t\cdot \dot{\bm \omega}^t - \dot{\bm \omega}^t \cdot {\bm \tau}^t)
\]
and then 
\[
{\bm \tau}^{t+\delta t} 
= 2 \eta_{\rm ve}\dot{\bm \varepsilon}^{d,t+\delta t} 
+ \frac{\eta_{\rm ve} }{\mu\delta t}  {\bm\tau}^t
- \frac{\eta_{\rm ve} }{\mu} ( {\bm \tau}^t\cdot \dot{\bm \omega}^t - \dot{\bm \omega}^t \cdot {\bm \tau}^t)
\]

We can further work out
\begin{eqnarray}
\frac{\eta_{\rm ve} }{ \mu\delta t} 
&=& \frac{\eta\mu \; \delta t}{\eta + \mu \delta t} \frac{1 }{ \mu\delta t} 
= \frac{\eta }{\eta + \mu \delta t} 
= \frac{1 }{1 + \frac{\mu}{\eta} \delta t}
\nn\\ 
\frac{\eta_{\rm ve} }{\mu}
&=&
\frac{\eta\mu \; \delta t}{\eta + \mu \delta t} \frac{1 }{\mu} 
= \frac{\eta \; \delta t}{\eta + \mu \delta t} 
\end{eqnarray}
so that finally
\[
{\bm \tau}^{t+\delta t} = ...
\]
We find this equation in \cite{modm02}
{\color{red} finish!}

$\eta/\mu$ is called the Maxwell relation time (also sometimes called the shear relaxation time \cite{modm02}).




